\documentclass[10pt, aspectratio=169]{beamer}

\usepackage{iftex}
\ifPDFTeX
  \usepackage[utf8]{inputenc}
  \usepackage[T1]{fontenc}
  \usepackage{lmodern}
\fi
\usepackage[brazilian]{babel}
\usepackage{graphicx}
\usepackage{booktabs}

% Cores Rangoo / FUMP-UFMG
\definecolor{rangoogreen}{HTML}{1A5C4A}

\usetheme[
  accent=rangoogreen,
  progressbar=foot,
  sectionpage=progressbar,
  subsectionpage=none,
  numbering=fraction,
  numberrail=section,
  block=fill,
  mathfonts=none,
  closing=contact
]{cookie}

% Metadados
\title{Rangoo}
\subtitle{Recarga de crédito do RU via PIX --- Apresentação Final}
\author[Boliñas de Gorfe]{%
  \begin{tabular}{@{}l@{}}
    Igor de Oliveira Martins dos Santos \\
    Ítalo Leal Lana Santos \\
    Vitor Hugo Dias Santos
  \end{tabular}%
}
\cookieuni{Universidade Federal de Minas Gerais (UFMG)}
\cookietagline{Hackathon InovaRU 2026/01 --- FUMP}
\date{Belo Horizonte, MG}

\cookieclosingtitle{Obrigado!}
\cookieclosingtext{Perguntas ou sugestões?}
\cookieclosingqr{https://github.com/igormartins4/app-hackathon-inova-ru}
\cookieclosingqrlabel{Código e slides no GitHub}

\begin{document}

\maketitle

\section{A Jornada de Gerson}

\begin{frame}{Gerson antes do Rangoo: Jantar contra o Relógio}
  \begin{columns}[T,onlytextwidth]
    \column{.48\textwidth}
      \begin{block}{O Tempo Limite}
        \small
        \begin{itemize}
          \item Aula acaba às 18h40 e a próxima inicia às 19h00.
          \item Apenas 20 minutos de janela para jantarem no RU.
          \item Pressão de tempo e correria diária.
        \end{itemize}
      \end{block}

    \column{.48\textwidth}
      \begin{block}{Sem Dinheiro Físico}
        \small
        \begin{itemize}
          \item Se tivesse dinheiro em papel, passaria rápido.
          \item Como não tem, é obrigado a encarar a enorme fila de pagamento com cartão de débito/crédito.
        \end{itemize}
      \end{block}
  \end{columns}
\end{frame}

\begin{frame}{Gerson antes do Rangoo: O Gargalo do Caixa Físico}
  \begin{columns}[T,onlytextwidth]
    \column{.48\textwidth}
      \begin{alertblock}{O Processo Manual no Caixa}
        \small
        \begin{itemize}
          \item Fila do caixa anda muito lentamente devido ao fluxo manual.
          \item O caixa precisa pegar a carteirinha física, verificar o saldo, digitar o valor na máquina de cartão, aguardar o aluno passar/digitar a senha, esperar a aprovação e só então liberar.
        \end{itemize}
      \end{alertblock}

    \column{.48\textwidth}
      \begin{block}{Consequência Acadêmica}
        \small
        \begin{itemize}
          \item Tempo gasto na fila do caixa compromete a janta rápida.
          \item Gerson perde o início da segunda aula (19h00).
        \end{itemize}
      \end{block}
  \end{columns}
\end{frame}

\begin{frame}{Gerson com o Rangoo: Pagamento Descentralizado}
  \begin{columns}[T,onlytextwidth]
    \column{.48\textwidth}
      \begin{block}{Zero Filas de Caixa}
        \small
        \begin{itemize}
          \item Não há mais necessidade de atendimento manual ou máquinas de cartão na entrada.
          \item Processo descentralizado: o estudante é dono do seu fluxo.
        \end{itemize}
      \end{block}

    \column{.48\textwidth}
      \begin{exampleblock}{Recarga Autônoma}
        \small
        \begin{itemize}
          \item Recarga rápida via PIX feita antecipadamente ou na própria fila de acesso.
          \item Saldo atualizado em segundos através de polling resiliente.
        \end{itemize}
      \end{exampleblock}
  \end{columns}
\end{frame}

\begin{frame}{Gerson com o Rangoo: A Catraca Inteligente}
  \begin{columns}[T,onlytextwidth]
    \column{.48\textwidth}
      \begin{block}{Acesso Expresso}
        \small
        \begin{itemize}
          \item Com saldo na conta, Gerson apenas aproxima a carteirinha/celular no sensor da catraca.
          \item Funcionamento semelhante a um validador de ônibus: débito automático instantâneo e acesso liberado.
        \end{itemize}
      \end{block}

    \column{.48\textwidth}
      \begin{block}{Resiliência e Rotina}
        \small
        \begin{itemize}
          \item Acesso rápido garante o jantar em poucos minutos.
          \item Gerson entra sem atrasos na aula das 19h00.
          \item Caching offline permite consultar saldo/cardápio mesmo sem sinal de rede móvel no RU.
        \end{itemize}
      \end{block}
  \end{columns}
\end{frame}

\section{O Produto Entregue}

\begin{frame}{Início, Saldo e Restaurantes}
  \begin{columns}[T,onlytextwidth]
    \column{.48\textwidth}
      \begin{block}{Dashboard Principal}
        \small
        \begin{itemize}
          \item \textbf{Saldo em Destaque}: Visão imediata do valor disponível.
          \item \textbf{Alerta de Saldo Baixo}: Banner inteligente para incentivar a recarga preventiva.
          \item \textbf{Carousel de Avisos}: Informes administrativos importantes da FUMP exibidos na Home.
        \end{itemize}
      \end{block}

    \column{.48\textwidth}
      \begin{exampleblock}{Contexto e Favoritos}
        \small
        \begin{itemize}
          \item \textbf{Favoritos com Status}: Lista rápida de RUs com horários e nível de fila.
          \item \textbf{Dados do Consumidor}: Detalhes sobre tipo de consumidor, centro de custo e situação (ativo/bloqueado).
        \end{itemize}
      \end{exampleblock}
  \end{columns}
\end{frame}

\begin{frame}{Recarga Expressa via PIX}
  \begin{columns}[T,onlytextwidth]
    \column{.48\textwidth}
      \begin{block}{Fácil Seleção de Valores}
        \small
        \begin{itemize}
          \item Chips de acesso rápido (R\$ 5, 10, 15, 20) ou entrada livre.
          \item Validação inteligente contra o limite de saldo de recarga individual (Zod).
        \end{itemize}
      \end{block}

    \column{.48\textwidth}
      \begin{alertblock}{QR Code e Confirmação}
        \small
        \begin{itemize}
          \item Código Pix Copia-e-Cola e exibição de QR Code Base64.
          \item Timer de expiração (2 min) e Polling com backoff exponencial.
          \item Confirmação de recebimento com saldo atualizado e comprovante.
        \end{itemize}
      \end{alertblock}
  \end{columns}
\end{frame}

\begin{frame}{Cardápio, Restaurantes e Histórico}
  \begin{columns}[T,onlytextwidth]
    \column{.48\textwidth}
      \begin{block}{Cardápio e Informações de RUs}
        \small
        \begin{itemize}
          \item Calendário semanal de cardápios e tags de restrição alimentar.
          \item Horários de funcionamento e rotas integradas via Google Maps.
          \item Scraping resiliente do cardápio público da FUMP.
        \end{itemize}
      \end{block}

    \column{.48\textwidth}
      \begin{exampleblock}{Histórico Unificado}
        \small
        \begin{itemize}
          \item Abas divididas entre refeições consumidas e recargas realizadas.
          \item Filtros por data de consulta.
          \item Scroll infinito com paginação automática baseada no contrato oficial.
        \end{itemize}
      \end{exampleblock}
  \end{columns}
\end{frame}

\begin{frame}{Transferência e Customizações do App}
  \begin{columns}[T,onlytextwidth]
    \column{.48\textwidth}
      \begin{block}{Transferência de Saldo (Bônus)}
        \small
        \begin{itemize}
          \item Envio de créditos de forma rápida entre estudantes através do CPF.
          \item Geração de comprovante compartilhável nativamente no dispositivo.
        \end{itemize}
      \end{block}

    \column{.48\textwidth}
      \begin{exampleblock}{Visual e Temas}
        \small
        \begin{itemize}
          \item Temas Escuro, Claro e Alto Contraste (Claro/Escuro) para ambientes externos.
          \item Respeito automático ao modo do sistema ou seleção manual no Perfil.
        \end{itemize}
      \end{exampleblock}
  \end{columns}
\end{frame}

\section{Engenharia e Segurança}

\begin{frame}{Segurança de Dados e APIs Oficiais}
  \begin{columns}[T,onlytextwidth]
    \column{.52\textwidth}
      \begin{block}{Segurança dos Dados}
        \small
        \begin{itemize}
          \item \textbf{Token JWT}: Criptografado via \texttt{expo-secure-store} (hardware - Android Keystore).
          \item \textbf{Sem Senha Salva}: Senha institucional em memória volátil, nunca gravada em disco.
          \item \textbf{Privacidade}: CPF omitido de requisições e URLs (obtido do JWT no backend).
        \end{itemize}
      \end{block}

    \column{.44\textwidth}
      \begin{exampleblock}{Integração com a FUMP}
        \small
        \begin{itemize}
          \item Comunicação via HTTPS obrigatório (TLS 1.2+).
          \item Endpoints alinhados ao contrato InovaRU v2.0.
          \item Validações dinâmicas contra o limite e teto (R\$ 5 a R\$ 500).
        \end{itemize}
      \end{exampleblock}
  \end{columns}
\end{frame}

\begin{frame}{Resiliência em Rede Instável (Offline)}
  \begin{columns}[T,onlytextwidth]
    \column{.52\textwidth}
      \begin{block}{Estratégia Offline-First}
        \small
        \begin{itemize}
          \item \textbf{TanStack Query}: Exibe cache imediatamente (SWR) e revalida dados em background.
          \item \textbf{AsyncStorage Cache}: Persiste cardápios, horários e histórico para uso sem sinal.
          \item \textbf{Offline Banner}: Notificação visual discreta, sem travar a navegação.
        \end{itemize}
      \end{block}

    \column{.44\textwidth}
      \begin{alertblock}{Algoritmo de Polling do PIX}
        \small
        \begin{itemize}
          \item \textbf{Backoff Exponencial}: Intervalos dinâmicos crescentes: 3s $\to$ 5s $\to$ 8s $\to$ 13s.
          \item \textbf{Jitter Aleatório}: Ruído de $\pm$1s para evitar picos concorrentes na API.
          \item \textbf{Timeout Seguro}: Limite de 2 minutos para poupar bateria e plano de dados.
        \end{itemize}
      \end{alertblock}
  \end{columns}
\end{frame}

\begin{frame}{Tratamento Estrito de Erros (Critério do Edital)}
  \begin{columns}[T,onlytextwidth]
    \column{.50\textwidth}
      \begin{block}{Mensagens Centralizadas}
        \small
        Erros de rede e HTTP mapeados em mensagens amigáveis:
        \begin{itemize}
          \item \textbf{401}: "Usuário ou senha inválidos."
          \item \textbf{404}: "Conta inativa. Procure a FUMP."
          \item \textbf{422}: "Valor fora do limite permitido."
          \item \textbf{429}: "Muitas tentativas." (Lê header \texttt{Retry-After}).
          \item \textbf{500}: "Ocorreu um erro. Tente novamente."
        \end{itemize}
      \end{block}

    \column{.46\textwidth}
      \begin{exampleblock}{Painel de Cenários Demo}
        \small
        Seletor no Perfil para testes rápidos de comportamento:
        \begin{itemize}
          \item Conta bloqueada ou inativa (Situação B).
          \item Pix pendente, aprovado, rejeitado ou expirado.
          \item Retornos HTTP específicos (401, 429, 500) e rede offline.
        \end{itemize}
      \end{exampleblock}
  \end{columns}
\end{frame}

\section{Acessibilidade e Inclusão}

\begin{frame}{Acessibilidade de Ponta a Ponta (WCAG)}
  \begin{columns}[T,onlytextwidth]
    \column{.50\textwidth}
      \begin{block}{Interface para Todo Mundo}
        \small
        \begin{itemize}
          \item \textbf{Leitores de Tela}: TalkBack e VoiceOver com rótulos descritivos (\texttt{accessibilityLabel}).
          \item \textbf{Área de Toque}: Targets com alvo mínimo de $48 \times 48\text{dp}$.
          \item \textbf{Fontes Escaláveis}: \texttt{ScaledText} adapta fontSize e lineHeight do sistema operacional.
        \end{itemize}
      \end{block}

    \column{.46\textwidth}
      \begin{alertblock}{Além das Regras do Edital}
        \small
        \begin{itemize}
          \item \textbf{Contraste WCAG AAA (7:1)}: Em saldo e valores para leitura sob luz solar.
          \item \textbf{Feedback Redundante}: Estados sinalizados com cor, texto e ícone.
          \item \textbf{Redução de Movimento}: Listener de preferência do sistema em tempo real.
        \end{itemize}
      \end{alertblock}
  \end{columns}
\end{frame}

\section{Fechamento}

\begin{frame}{Stack Tecnológica e Diferenciais}
  \begin{columns}[T,onlytextwidth]
    \column{.52\textwidth}
      \begin{block}{A Stack do Rangoo}
        \small
        \begin{tabular}{@{}ll@{}}
          \toprule
          \textbf{Camada} & \textbf{Tecnologia} \\
          \midrule
          Mobile & React Native (Expo 57) \\
          Navegação & Expo Router 6 \\
          Estado (UI) & Zustand 5 \\
          Estado (Server) & TanStack Query 5 \\
          Estilos & NativeWind v4 \\
          Segurança & Expo Secure Store \\
          \bottomrule
        \end{tabular}
      \end{block}

    \column{.44\textwidth}
      \begin{exampleblock}{Diferenciais de Entrega}
        \small
        \begin{itemize}
          \item \textbf{Software Livre}: Licença MIT garantindo conformidade legal com a FUMP.
          \item \textbf{Zero Fricção de Login}: Uso de credenciais institucionais já existentes.
          \item \textbf{Arquitetura Modular}: Organizado por features separadas, facilitando a manutenção e escala.
        \end{itemize}
      \end{exampleblock}
  \end{columns}
\end{frame}

\begin{frame}{Demonstração Prática (Live Demo)}
  \begin{block}{Apresentação em Tempo Real}
    \small
    Faremos agora a demonstração prática do Rangoo Universitário diretamente na tela do celular:
    \begin{enumerate}
      \item \textbf{Login e Dashboard}: Acesso rápido, visualização de saldo e aviso inteligente.
      \item \textbf{Recarga Expressa}: PIX Copia-e-Cola e confirmação instantânea em menos de 30 segundos.
      \item \textbf{Cardápio \& Histórico}: Consulta às opções do dia e extrato financeiro/consumo.
      \item \textbf{Painel de Testes}: Simulação de erro 429 (Rate Limit), conta inativa (404) e rede offline.
      \item \textbf{Transferência de Créditos}: Envio imediato de saldo para outro estudante.
    \end{enumerate}
  \end{block}
\end{frame}

\end{document}
